% ============================================================
%  Appendix B --- Interview Preparation Guide
% ============================================================
\chapter{Interview Preparation Guide}
\label{ch:interview-prep}

This appendix consolidates all interview-relevant questions and answers from the preceding chapters, organised by topic.  Each answer references your \shiftdb{} implementation so you can speak from first-hand experience.

% ═══════════════════════════════════════════════════
\section{Storage Engine \& Page Management}
% ═══════════════════════════════════════════════════

\begin{interviewbox}[Q1: What is a slotted-page and why is it used?]
A slotted-page stores a header, a slot directory (growing forward), and tuple data (growing backward).  Each slot records the offset and length of a tuple within the page.  This layout supports \textbf{variable-length records} and allows tuples to be moved during compaction --- only the slot entry needs updating, not external pointers.

\emph{In \shiftdb{}:} See \code{Page} class in \code{page.h}.  The header is 32 bytes (\code{PageHeader}); each slot is 12 bytes (\code{SlotEntry}).
\end{interviewbox}

\begin{interviewbox}[Q2: How does \shiftdb{} handle NULL values in tuple storage?]
A \textbf{null bitmap} at the start of each serialised row uses $\lceil n/8 \rceil$ bytes (one bit per column).  If a bit is set, the corresponding field is skipped entirely in the byte stream.  This is the same approach used by PostgreSQL and MySQL.

\emph{Code ref:} \code{Table::serializeRow()} in \code{table.h}.
\end{interviewbox}

\begin{interviewbox}[Q3: What happens when all pages are full during an INSERT?]
\code{Table::insertRow()} iterates through all existing pages.  If none have sufficient free space, it calls \code{DiskManager::allocatePage()} to append a fresh 4\,KB page to the \code{.nxdb} file and inserts the tuple there.
\end{interviewbox}

\begin{interviewbox}[Q4: Compare heap file vs.\ sorted file vs.\ LSM tree.]
\begin{itemize}[nosep]
  \item \textbf{Heap file} (unordered): $\bigO{1}$ insert, $\bigO{n}$ search.  Simple, good for OLTP with indexes.
  \item \textbf{Sorted file}: $\bigO{n}$ insert (must shift), $\bigO{\log n}$ search.  Good for range queries.
  \item \textbf{LSM tree}: $\bigO{1}$ insert (write to memtable), $\bigO{\log n}$ search (check multiple levels).  Optimised for write-heavy workloads (RocksDB, Cassandra).
\end{itemize}
\emph{\shiftdb{} uses a heap file} for its simplicity.
\end{interviewbox}

% ═══════════════════════════════════════════════════
\section{Buffer Management}
% ═══════════════════════════════════════════════════

\begin{interviewbox}[Q5: What is a buffer pool and why does it matter?]
The buffer pool is an in-memory cache of disk pages.  It exploits temporal locality: recently accessed pages are likely to be accessed again.  A well-tuned buffer pool serves 95\%+ of requests from memory, avoiding slow disk I/O.

\emph{In \shiftdb{}:} \code{BufferPool} class, default size 1024 pages (4\,MB).
\end{interviewbox}

\begin{interviewbox}[Q6: Explain LRU eviction and its weakness.]
LRU evicts the least recently used unpinned page.  Its weakness is \textbf{sequential flooding}: a full table scan touches every page once, evicting frequently-used pages from other queries.

\textbf{Fix:} LRU-K (track $K$ most recent access times), 2Q (separate queues for hot and cold pages), or a dedicated ring buffer for sequential scans (PostgreSQL).
\end{interviewbox}

\begin{interviewbox}[Q7: What is a dirty page and when is it flushed?]
A dirty page has been modified in memory but not yet written to disk.  It is flushed: (1) on eviction (before the frame is reused), (2) during a checkpoint, (3) at clean shutdown.
\end{interviewbox}

\begin{interviewbox}[Q8: What is pin count?]
Pin count tracks how many active operations hold a reference to a page.  A page with \code{pinCount > 0} \textbf{cannot be evicted}.  Callers must call \code{unpinPage()} when done.
\end{interviewbox}

% ═══════════════════════════════════════════════════
\section{Indexing}
% ═══════════════════════════════════════════════════

\begin{interviewbox}[Q9: Why do databases use B$^+$-trees instead of B-trees?]
B$^+$-trees store data only in leaf nodes, so internal nodes hold more keys per page (lower tree height).  Leaf nodes are linked in a doubly-linked list, making range scans a simple sequential traversal instead of requiring tree re-entry.
\end{interviewbox}

\begin{interviewbox}[Q10: What is a clustered index?]
A clustered index determines the \textbf{physical order} of rows on disk.  A table can have at most one.  In MySQL/InnoDB, the primary key is always the clustered index.  Non-clustered (secondary) indexes store pointers back to the primary key.
\end{interviewbox}

\begin{interviewbox}[Q11: When would you choose a hash index?]
For \textbf{exact-match} queries only (\code{WHERE id = 42}).  Hash indexes provide $\bigO{1}$ amortised lookups but cannot serve range scans, \code{ORDER BY}, or prefix \code{LIKE} queries.  B$^+$-trees are the default because they handle all these.
\end{interviewbox}

\begin{interviewbox}[Q12: What triggers a B-tree node split?]
When a node exceeds $m-1$ keys (for an order-$m$ tree).  The median key is pushed up to the parent.  If the parent also overflows, the split cascades upward.  In the worst case, the root splits and tree height increases by one.
\end{interviewbox}

% ═══════════════════════════════════════════════════
\section{SQL Parsing}
% ═══════════════════════════════════════════════════

\begin{interviewbox}[Q13: What is a recursive-descent parser?]
An RDP has one function per grammar non-terminal.  Each function matches the current tokens against its production rules.  It is top-down, predictive, and easy to write by hand.

\textbf{Limitation:} Cannot handle left-recursive grammars directly.  \shiftdb{} handles operator precedence by nesting calls: \code{parseOr} $\to$ \code{parseAnd} $\to$ \code{parseComparison} $\to$ \code{parseAddition} $\to$ etc.
\end{interviewbox}

\begin{interviewbox}[Q14: What is the time complexity of tokenisation and parsing?]
Both are $\bigO{n}$ where $n$ is the length of the SQL string.  The tokeniser makes one pass; the parser makes one pass over the token list with no backtracking.
\end{interviewbox}

% ═══════════════════════════════════════════════════
\section{Query Optimisation}
% ═══════════════════════════════════════════════════

\begin{interviewbox}[Q15: Logical plan vs.\ physical plan?]
A \textbf{logical plan} specifies \emph{what} operations to perform (Select, Project, Join) using relational algebra.  A \textbf{physical plan} specifies \emph{how}: which algorithm (hash join, sort-merge), which indexes, scan direction, etc.
\end{interviewbox}

\begin{interviewbox}[Q16: Explain predicate pushdown.]
Move filter conditions as close to the data source as possible.  This reduces the number of rows flowing through expensive operators (especially joins).  Example: push \code{WHERE city = 'NYC'} below a join so that only NYC customers are joined.
\end{interviewbox}

\begin{interviewbox}[Q17: Compare nested loop, hash join, and sort-merge join.]
\begin{itemize}[nosep]
  \item \textbf{Nested loop}: $\bigO{n \times m}$.  Simple; good for small tables or with an index on the inner table.
  \item \textbf{Hash join}: $\bigO{n + m}$.  Builds a hash table on the smaller side.  Best for equality joins with sufficient memory.
  \item \textbf{Sort-merge join}: $\bigO{n\log n + m\log m}$.  Sorts both sides, then merges.  Good when inputs are already sorted or when the result needs ordering.
\end{itemize}
\end{interviewbox}

\begin{interviewbox}[Q18: Volcano vs.\ materialisation vs.\ vectorised execution?]
\begin{itemize}[nosep]
  \item \textbf{Volcano (iterator)}: one tuple at a time via \code{next()}.  Low memory, high per-tuple overhead.
  \item \textbf{Materialisation}: full intermediate results.  Simple but memory-heavy.  \emph{(Used by \shiftdb{}.)}
  \item \textbf{Vectorised}: batches of tuples per \code{next()}.  Amortises function-call overhead, cache-friendly.  (DuckDB.)
\end{itemize}
\end{interviewbox}

% ═══════════════════════════════════════════════════
\section{Transactions \& Concurrency}
% ═══════════════════════════════════════════════════

\begin{interviewbox}[Q19: Explain ACID with a real-world example.]
Bank transfer: \emph{Atomicity} = debit and credit both happen or neither.  \emph{Consistency} = total balance is preserved.  \emph{Isolation} = concurrent transfers do not interfere.  \emph{Durability} = committed transfer survives a crash.
\end{interviewbox}

\begin{interviewbox}[Q20: What is Two-Phase Locking (2PL)?]
In 2PL, a transaction acquires locks in a \emph{growing phase} and releases them in a \emph{shrinking phase}.  No lock is acquired after any lock is released.  \textbf{Strict 2PL} holds all locks until commit, preventing cascading aborts.
\end{interviewbox}

\begin{interviewbox}[Q21: What are the SQL isolation levels?]
\begin{enumerate}[nosep]
  \item \textbf{Read Uncommitted}: see uncommitted data (dirty reads allowed).
  \item \textbf{Read Committed}: only see committed data (default in PostgreSQL, Oracle).
  \item \textbf{Repeatable Read}: re-reading the same row gives the same result (default in MySQL/InnoDB).
  \item \textbf{Serializable}: transactions behave as if executed serially.
\end{enumerate}
\end{interviewbox}

\begin{interviewbox}[Q22: How do you detect and resolve deadlocks?]
\textbf{Detection:} Build a wait-for graph; if a cycle exists, there is a deadlock.  \textbf{Resolution:} Abort the youngest (or cheapest) transaction in the cycle.  \textbf{Prevention:} wound-wait or wait-die schemes use timestamps to break ties.
\end{interviewbox}

% ═══════════════════════════════════════════════════
\section{MVCC}
% ═══════════════════════════════════════════════════

\begin{interviewbox}[Q23: What is MVCC and why is it used?]
MVCC keeps multiple timestamped versions of each row.  Readers see a consistent snapshot without acquiring locks, eliminating reader-writer contention.  Used by PostgreSQL, MySQL/InnoDB, Oracle, SQL Server.
\end{interviewbox}

\begin{interviewbox}[Q24: PostgreSQL MVCC vs.\ Oracle MVCC?]
\textbf{PostgreSQL:} Stores all versions in the same heap table; dead versions cleaned by VACUUM.  Can cause table bloat if VACUUM falls behind.

\textbf{Oracle:} Main table stores latest version; old versions in undo segments (circular buffer).  No bloat, but limited read-back window (``snapshot too old'' error).
\end{interviewbox}

\begin{interviewbox}[Q25: What is the write skew anomaly?]
Under snapshot isolation, two transactions read overlapping data and make disjoint writes that together violate a constraint.  Example: two doctors both check that the other is on call and remove themselves --- result: no one on call.

Fix: use Serializable isolation (predicate locking / SSI).
\end{interviewbox}

% ═══════════════════════════════════════════════════
\section{Write-Ahead Logging \& Recovery}
% ═══════════════════════════════════════════════════

\begin{interviewbox}[Q26: What is the WAL principle?]
\textbf{``Log before data.''} Before any modified page is written to disk, all log records for that page must be flushed to the log file.  This guarantees enough information exists to redo or undo any change after a crash.
\end{interviewbox}

\begin{interviewbox}[Q27: Describe the three phases of ARIES.]
\begin{enumerate}[nosep]
  \item \textbf{Analysis:} Scan log from last checkpoint; build Dirty Page Table and Active Transaction Table.
  \item \textbf{Redo:} Replay all changes from the earliest recLSN forward (even for aborted transactions).
  \item \textbf{Undo:} Roll back all uncommitted transactions in reverse LSN order; write CLRs for idempotency.
\end{enumerate}
\end{interviewbox}

\begin{interviewbox}[Q28: What is STEAL/NO-FORCE?]
\textbf{STEAL:} Buffer pool can flush uncommitted dirty pages (needs undo on crash).  \textbf{NO-FORCE:} Committed data is not forced to disk at commit (needs redo on crash).  ARIES uses STEAL/NO-FORCE for maximum flexibility and performance.
\end{interviewbox}

% ═══════════════════════════════════════════════════
\section{Architecture \& Design}
% ═══════════════════════════════════════════════════

\begin{interviewbox}[Q29: Describe the architecture of a DBMS.]
Five core components: (1) \textbf{Query processor} (parser, optimiser, executor); (2) \textbf{Storage engine} (page management, file I/O); (3) \textbf{Buffer manager} (page cache); (4) \textbf{Catalog} (metadata); (5) \textbf{Transaction/concurrency manager}.

\emph{\shiftdb{} implements all five.}  The presentation layer (React UI) and network layer (HTTP server) sit on top.
\end{interviewbox}

\begin{interviewbox}[Q30: Why is layer separation important?]
\begin{itemize}[nosep]
  \item \textbf{Modularity:} each component can evolve independently.
  \item \textbf{Testability:} layers can be unit-tested in isolation.
  \item \textbf{Portability:} the storage layer can be swapped (file system to cloud storage) without touching the parser.
\end{itemize}
\end{interviewbox}

% ═══════════════════════════════════════════════════
\section{Performance \& Optimisation}
% ═══════════════════════════════════════════════════

\begin{interviewbox}[Q31: How would you improve full table scan performance?]
(1) Add B$^+$-tree indexes for point/range queries.  (2) Use column projection to skip unneeded columns.  (3) Use a ring buffer for sequential scans to prevent cache pollution.  (4) Parallelise across threads.  (5) Vectorised tuple processing.
\end{interviewbox}

\begin{interviewbox}[Q32: What metrics would you track for slow queries?]
Buffer pool hit rate, query execution time, pages read (I/O), rows scanned vs.\ returned (selectivity), number and type of joins, CPU time, lock waits, and disk I/O latency.
\end{interviewbox}

% ═══════════════════════════════════════════════════
\section{Web \& Systems Design}
% ═══════════════════════════════════════════════════

\begin{interviewbox}[Q33: Why build a custom HTTP server?]
Zero external dependencies --- the entire C++ backend compiles without third-party libraries (except nlohmann/json for JSON handling).  This simplifies the build and demonstrates low-level networking.  In production, use a battle-tested library for robustness and security.
\end{interviewbox}

\begin{interviewbox}[Q34: What is CORS?]
Cross-Origin Resource Sharing.  Browsers block HTTP requests from one origin (e.g.\ \code{localhost:5173}) to another (e.g.\ \code{localhost:8080}) unless the server includes \code{Access-Control-Allow-Origin} headers.  \shiftdb{}'s server sends \code{*} to allow all origins.
\end{interviewbox}

\begin{interviewbox}[Q35: How would you scale \shiftdb{} to handle production traffic?]
\begin{enumerate}[nosep]
  \item \textbf{Thread pool} instead of one thread per connection.
  \item \textbf{WAL} for durability and crash recovery.
  \item \textbf{Index-accelerated queries} for sub-millisecond lookups.
  \item \textbf{Connection pooling} on the frontend.
  \item \textbf{Replication} for read scaling (leader-follower).
  \item \textbf{Sharding} for write scaling (hash or range partitioning).
\end{enumerate}
\end{interviewbox}

% ═══════════════════════════════════════════════════
\section{Rapid-Fire Questions}
% ═══════════════════════════════════════════════════

\begin{longtable}{p{0.55\textwidth} p{0.40\textwidth}}
\toprule
\textbf{Question} & \textbf{One-Line Answer} \\
\midrule
\endhead
What data structure does InnoDB use for its primary index? & B$^+$-tree (clustered) \\
Default isolation level in PostgreSQL? & Read Committed \\
Default isolation level in MySQL/InnoDB? & Repeatable Read \\
What does EXPLAIN do? & Shows the query execution plan \\
Difference between DELETE and TRUNCATE? & DELETE logs each row (can rollback); TRUNCATE deallocates pages \\
What is a covering index? & Index that contains all columns needed by the query \\
What is index-only scan? & Query answered entirely from the index without reading the heap \\
What is a composite index? & Index on multiple columns, e.g.\ \code{(last\_name, first\_name)} \\
Leftmost prefix rule? & A composite index on \code{(A,B,C)} can serve queries on \code{A}, \code{A,B}, or \code{A,B,C} \\
What is a correlated subquery? & A subquery that references columns from the outer query \\
What is normalisation? & Organising tables to reduce redundancy (1NF, 2NF, 3NF, BCNF) \\
What is denormalisation? & Intentionally adding redundancy for read performance \\
What is a view? & A stored query that acts as a virtual table \\
What is a materialised view? & A view whose results are cached on disk \\
CAP theorem? & Consistency, Availability, Partition tolerance --- pick two \\
\bottomrule
\end{longtable}
